\documentclass[10pt, aspectratio=169]{beamer}

\usepackage[utf8]{inputenc}
\usepackage[T1]{fontenc}
\usepackage[ngerman,shorthands=off]{babel}
\usepackage{graphicx}
\usepackage[absolute,overlay]{textpos}
\usetheme{Madrid}
\usecolortheme{seahorse} 
\setbeamertemplate{navigation symbols}{}
\usepackage[backend=biber]{biblatex}
\addbibresource{/Users/jakubmarczak/Downloads/literatur.bib}
\usepackage{csquotes}
\usepackage{hyperref}
\usepackage{wasysym}
\usepackage{xcolor}
\usepackage{tikz}
\definecolor{progress}{HTML}{9b9bba}

\title{Projekt \MakeUppercase{\romannumeral 2}  Zwischenpräsentation}
\author[J. Marczak, G. Kathiravan, M. I. Rahman]{Jakub Marczak, Ganeisraaj Kathiravan, Muhammad Ibrahim Rahman}

\setbeamertemplate{footline}{
  \leavevmode
  
  % Infoleiste
  \hbox{
    \begin{beamercolorbox}[wd=.333333\paperwidth,ht=2.5ex,dp=1ex,left]{author in head/foot}
      \hspace{1em}\insertshortauthor
    \end{beamercolorbox}
    \begin{beamercolorbox}[wd=.333333\paperwidth,ht=2.5ex,dp=1ex,center]{title in head/foot}
      \insertshorttitle
    \end{beamercolorbox}
    \begin{beamercolorbox}[wd=.333333\paperwidth,ht=2.5ex,dp=1ex,right]{date in head/foot}
      \insertframenumber{} / \inserttotalframenumber \hspace{1em}
    \end{beamercolorbox}
  }

  % Progressbar
  \hbox{
    \begin{beamercolorbox}[wd=\paperwidth,ht=1.3ex,dp=0ex,left]{}
      \tikz{
        \fill[gray!15] (0,0) rectangle (\paperwidth,1.3ex);
        \fill[progress] (0,0) rectangle
        ({\paperwidth*\insertframenumber/\inserttotalframenumber},1.3ex);
      }
    \end{beamercolorbox}
  }
}

\begin{document}

\begin{frame}[plain]
  \titlepage
\end{frame}

% Folie 2
\begin{frame}
    \frametitle{Inhalt}
    \begin{enumerate}
        \item Population verstehen
        \item Nierenfunktion unter Vancomycin-Therapie
        \begin{itemize} 
        \item[$\rightarrow$] Verhalten von SCr- und eGFR-Werten
        \item[$\rightarrow$] Auswirkung vom Vancomycin-Spiegel auf Nierenmarker
        \item[$\rightarrow$] Zusammenhang zwischen Vancomycin-Spiegel und Nierenfunktion
        \item[$\rightarrow$] Einfluss der Nierenfunktion auf Gesundheits-Marker
      \end{itemize}
        \item Mortalitätsanalyse
        \item[$\rightarrow$] Welche Einflüsse begleiten die Mortalität am meisten?
        \begin{itemize}
            \item Nierenfunktion
            \item Schweregrades der Erkrankung
            \item Komorbiditäten
            \item Nephrotoxine
        \end{itemize}
    \end{enumerate}
\end{frame}

% Folie 3 
\begin{frame}
    \frametitle{Population verstehen}
%    \begin{figure}
%        \centering
%        \includegraphics[width=1\linewidth]{population.pdf}
%        \vspace{-0.75cm}
%        \caption{Kerndichteschätzungen mit Gaußkern für Alter, Größe und Gewicht.}
%    \end{figure}
\end{frame}

% Folie 4 
\begin{frame}
    \frametitle{Verhalten von SCr- und eGFR-Werten}
%    \begin{figure}
%        \centering
%        \vspace{-0.2cm}
%        \includegraphics[width=0.6\linewidth]{corrplot nierenmarker.pdf}
%        \vspace{-0.4cm}
%        \caption{Korrelationsplot zwischen SCr- und eGFR-Werten als Nierenmarker.}
%    \end{figure}
\end{frame}

% Folie 5 
\begin{frame}
    \frametitle{Auswirkung vom Vancomycin-Spiegel auf Nierenmarker}
%    \begin{figure}
%        \centering
%        \includegraphics[width=1.02\linewidth]{ZhSCreGFR.pdf}
%        \vspace{-0.4cm}
%        \caption{Hexbin-Plots von $\Delta$SCr und $\Delta$eGFR gegen die mittlere Vancomycin-Konzentration ($\bar{C}$) mit \textit{LOESS}-Trendlinie.}
%    \end{figure}
\end{frame}

% Folie 6 
\begin{frame}
    \frametitle{Auswirkung vom Vancomycin-Spiegel auf Nierenmarker}
%    \begin{figure}
%        \centering
%        \includegraphics[width=1.02\linewidth]{HexScat.pdf}
%         \vspace{-0.4cm}
%        \caption{Vergleich der Darstellung desselben Datensatzes als Scatterplot und Hexbin-Plot.}
%    \end{figure}
\end{frame}

% Folie 7 
\begin{frame}
	\frametitle{Zusammenhang zwischen Vancomycin-Spiegel und Nierenfunktion}
%	\begin{figure}
%		\centering
%		\hspace*{1.5cm}
%		\includegraphics[width=0.7\linewidth]{C24xeGFRStart.pdf}
%		\vspace{-0.45cm} 
%		\caption{Zusammenhang zwischen Vancomycin-Serum-Spiegel nach 24h ($C_{24}$) und der geschätzten glomerulären Filtrationsrate zu Beginn der Therapie($eGFR_{Start}$) (GAM-Trend).}
%	\end{figure}
\end{frame}

% Folie 8 
\begin{frame}
	\frametitle{Einfluss der Nierenfunktion auf Gesundheits-Marker}
%	\begin{figure}
%		\centering
%		\includegraphics[width=0.82\linewidth]{SAPSSOFAeGFR.pdf}
%		\vspace{-0.4cm}
%		\caption{Korrelation von $eGFR_{Start}$ zum SAPS- und SOFA-Score mit LOESS-Trendlinie und 95 \% Konfidenzintervall.}
%	\end{figure}
\end{frame}

% Folie 9 
\begin{frame}
	\frametitle{Nierenfunktion in Verbindung mit Mortalität }
%	\begin{figure}
%		\centering
%		\includegraphics[width=0.88\linewidth]{eGFRxMort.pdf}
%		\vspace{-0.55cm}
%		\caption{Verteilung von $eGFR_{Start}$ nach Mortalitätsstatus.}
%	\end{figure}
\end{frame}

% Folie 10 
\begin{frame}
    \frametitle{Einfluss des Schweregrades der Erkrankung auf die Mortalität}
%    \begin{figure}
%        \centering
%        \includegraphics[width=0.65\linewidth]{SchwerxMort.pdf}
%         \vspace{-0.2cm}
%        \caption{Boxplots der Schweregradvariablen zu Therapiebeginn (erste 24h), gruppiert nach Mortalitätsstatus (lebt vs. gestorben).}
%    \end{figure}
\end{frame}

% Folie 11
\begin{frame}
    \frametitle{Welche Komorbiditäten begleiten die Mortalität am meisten?}
%    \begin{figure}
%        \centering
%        \includegraphics[width=1\linewidth]{KomorxMort.pdf}
%        \caption{Barplot Anzahl von Patienten und Mortalitätsrate abhängig von Komorbiditäten.}
%    \end{figure}
\end{frame}

% Folie 12
\begin{frame}
    \frametitle{Welche Nephrotoxine begleiten die Mortalität am meisten?}
%    \begin{figure}
%        \centering
%        \includegraphics[width=1\linewidth]{NephroxMort.pdf}
%        \caption{Barplot Anzahl von Patienten und Mortalitätsrate abhängig von Nephrotoxinen.}
%    \end{figure}
\end{frame}

% Folie 13 
\begin{frame}
    \frametitle{Therapiedauer nach Mortalitätsstatus}
%    \begin{figure}
%        \centering
%        \includegraphics[width=0.7\linewidth]{therapiedauer_mortalitaet.png}
%        \caption{Dauer der Vancomycin-Therapie nach Mortalitätsstatus.}
%    \end{figure}
\end{frame}

% Folie 14
\begin{frame}
    \frametitle{Mortalitätsrate nach C24-Quartilen}
%    \begin{figure}
%        \centering
%        \includegraphics[width=0.9\linewidth]{ Rplot.png}
%        \caption{Mortalitätsrate nach Quartilen des Vancomycin-Serumspiegels (C24).}
%    \end{figure}
\end{frame}

\end{document}
